\begin{Theorem}{Euler-Lagrange 方程}
\label{thm:elasticity-euler-lagrange-equations}
假设 $\overline{\bm{p}}\in\mathcal{D}_{\mathcal{E}}\cap C^2(\overline\Omega)$ 解 \eqref{eq:elasticity-equilibrium-global-minimization}, 且 \eqref{eq:elasticity-positive-orientation-determinant} 在 $\Omega$ 中几乎处处成立. 则在每个 $\bm{x}\in\Omega$ 处, $\overline{\bm{p}}$ 必满足 Euler-Lagrange 方程
\begin{equation}\label{eq:elasticity-euler-lagrange-equations}
-\sum_j\dfrac{\partial}{\partial x_j}\mathcal{A}_{ij}(\bm{x})+\dfrac{\partial\Psi}{\partial p_i}(\bm{p}(\bm{x}))=0,\quad i=1,2,3,
\end{equation}
其中
\begin{equation}\label{eq:elasticity-euler-lagrange-stress-coefficients}
\mathcal{A}_{ij}(\bm{x}):=\dfrac{\partial\widehat{\mathcal{W}}}{\partial F_{ij}}(\bm{x},\bm{F}(\bm{x})).
\end{equation}
在均匀材料中: $\mathcal{W}=\widehat{\mathcal{W}}(\bm{F})$, 有
\begin{equation}\label{eq:elasticity-homogeneous-stress-divergence}
\sum_j\dfrac{\partial}{\partial x_j}\mathcal{A}_{ij}(\bm{x})=\sum_{j,k,l}a_{ijkl}(\bm{x})\dfrac{\partial^2p_k}{\partial x_j\partial x_l},
\end{equation}
其中
\begin{equation}\label{eq:elasticity-homogeneous-second-derivative-coefficients}
a_{ijkl}(\bm{x}):=\dfrac{\partial^2\mathcal{W}}{\partial F_{ij}\partial F_{kl}}(\bm{F}(\bm{x})).
\end{equation}

\end{Theorem}

\begingroup
\renewcommand{\endProof}{\par\nobreak\hfill$\blacksquare$\par\vspace{0.1cm}}
\begin{Proof}
由于 $\overline{\bm{p}}\in\mathcal{D}_{\mathcal{E}}\cap C^2(\overline\Omega)$, 可知对任意 $\bm{\phi}=\sum\phi_i\bm{e}_i\in[\mathcal{D}(\Omega)]^3$, 当 $\epsilon$ 足够小时有
\begin{equation}\label{eq:elasticity-compact-variation-admissibility}
\overline{\bm{p}}+\epsilon\bm{\phi}\in\mathcal{D}_{\mathcal{E}}
\end{equation}
因此, 对任意 $\bm{\phi}$, 实值函数
\begin{equation}\label{eq:elasticity-scalar-variation-energy}
f(\epsilon):=\mathcal{E}(\overline{\bm{p}}+\epsilon\bm{\phi})
\end{equation}
在 $\epsilon=0$ 附近的某个区间内有良好定义. 此外, 由于假设 $\mathcal{E}$ 在 $\overline{\bm{p}}$ 处取极小值, $f$ 在 $\epsilon=0$ 处取极小值. 现在利用一致收敛的标准结果在积分号下求导, 可知 $f$ 可微. 得到
\begin{equation}\label{eq:elasticity-first-variation-integral}
\begin{aligned}
 f^{\prime}(\epsilon)=\int_\Omega&\sum_{i,j}\dfrac{\partial\mathcal{W}}{\partial F_{ij}}(\bm{x},\overline{\bm{F}}(\bm{x})+\epsilon\nabla\bm{\phi}(\bm{x}))\dfrac{\partial\phi_i}{\partial x_j}\\
 &+\sum_i\dfrac{\partial\Psi}{\partial p_i}(\overline{\bm{p}}+\epsilon\bm{\phi})\phi_i\,\mathrm{d}\bm{x}.
 \end{aligned}
\end{equation}
令导数在 $\epsilon=0$ 处等于零, 并作分部积分, 得到
\begin{equation}\label{eq:elasticity-first-variation-integrated-by-parts}
\int_\Omega\sum_i\left[-\sum_j\dfrac{\partial}{\partial x_j}\dfrac{\partial\mathcal{W}}{\partial F_{ij}}(\bm{x},\overline{\bm{F}}(\bm{x}))+\dfrac{\partial\Psi}{\partial p_i}(\overline{\bm{p}}(\bm{x}))\right]\phi_i(\bm{x})\,\mathrm{d}\bm{x}=0.
\end{equation}
由于 $\bm{\phi}$ 任意, 这说明 \eqref{eq:elasticity-euler-lagrange-equations} 在分布意义下成立. 由于方程中的每一项都连续, 它实际上必逐点成立.
\end{Proof}
\endgroup
